\documentclass[11pt]{article}
\usepackage[utf8]{inputenc}
\usepackage[T1]{fontenc}
\usepackage[italian]{babel}
\usepackage{geometry}
\geometry{a4paper, margin=1in}
\usepackage{xcolor}
\usepackage{listings}
\usepackage{hyperref}

\usepackage{tcolorbox}
\tcbuselibrary{breakable}
\begin{document}

\begin{tcolorbox}[colback=blue!5!white,colframe=blue!75!black,title=\textbf{DOMANDA}, breakable]
Esercizio 0\\
Effettuare una cattura Wireshark relativa al funzionamento della chat. Occorre attivare la cattura prima di aprire la URL dalla finestra del browser e ascoltare non solo sull’interfaccia di loopback ma anche su quella verso l’esterno (opzione ANY).
\end{tcolorbox}

\begin{tcolorbox}[colback=green!5!white,colframe=green!75!black,title=\textbf{RISPOSTA}, breakable]
\textbf{Soluzione Esercizio 0:} Cattura Wireshark del funzionamento della chat

\textbf{Operazioni Preliminari:}\\
Ho creato la cartella dell'esercizio (\texttt{es33 - WebSocket Esercizio 0}) e vi ho copiato dentro tutto il codice sorgente del progetto "11 WebSocket\_Chat".

\textbf{Analisi della cattura Wireshark:}\\
Avviando una cattura di rete sull'interfaccia "ANY" (o di loopback \texttt{lo} se testiamo tutto in locale) e filtrando per la porta del nostro server (es. \texttt{tcp.port == 4000}), aprendo il browser e mandando un messaggio in chat si assisterà a una precisa sequenza di eventi di rete che dimostra il funzionamento e il passaggio da \texttt{HTTP} a WebSocket.

Ecco cosa si vede passo per passo nei pacchetti:

\begin{enumerate}
    \item \textbf{FASE HTTP (Recupero File Statici)}
    \begin{itemize}
        \item Inizialmente si vedrà il classico "3-way handshake" del \texttt{TCP} (SYN, SYN-ACK, ACK) per stabilire la connessione con la porta 4000 del server Node.js.
        \item Il browser invia una \texttt{HTTP GET /} per richiedere la pagina principale.
        \item Il server Express risponde con un \texttt{HTTP/1.1 200 OK} che contiene il codice HTML della pagina.
        \item Subito dopo, il browser invierà altre richieste \texttt{HTTP GET} per scaricare i file statici referenziati nell'HTML (come \texttt{/chat.js} e \texttt{/styles.css}).
    \end{itemize}

    \item \textbf{FASE DI UPGRADE (Il passaggio a WebSocket)}
    \begin{itemize}
        \item Il client (grazie all'istruzione \texttt{io.connect()} nel file JavaScript) inizializzerà una richiesta speciale verso il server. Nello specifico, invierà una richiesta \texttt{HTTP GET} contenente degli header molto particolari:
        \begin{verbatim}
Connection: Upgrade
Upgrade: websocket
Sec-WebSocket-Key: [chiave casuale]
        \end{verbatim}
        \item Se il server supporta WebSocket (come nel nostro caso grazie a Socket.IO), risponderà al browser con un pacchetto speciale: \texttt{HTTP/1.1 101 Switching Protocols}.
        \item Negli header della risposta 101, il server confermerà l'avvenuto "cambio di lingua" ri-dichiarando \texttt{Upgrade: websocket} e rispondendo alla chiave di sicurezza.
    \end{itemize}

    \item \textbf{FASE WEBSOCKET (Comunicazione Bidirezionale in Real-Time)}
    \begin{itemize}
        \item Da questo momento in poi (dopo il pacchetto 101), Wireshark smetterà di classificare la conversazione come "HTTP". I pacchetti successivi sulla stessa connessione \texttt{TCP} saranno etichettati col protocollo "WebSocket".
        \item Quando l'utente preme il bottone per inviare un messaggio nella chat, Wireshark mostrerà un pacchetto "WebSocket Text Frame".
        \item Poiché la comunicazione non è cifrata (non stiamo usando WSS ma semplice WS), cliccando sul payload del pacchetto in Wireshark sarà possibile leggere in chiaro il messaggio inviato.
        \item Essendo mediato dalla libreria Socket.IO, il testo nel payload apparirà con un piccolo prefisso numerico usato dal protocollo interno della libreria, seguito dal nostro oggetto JSON. Ad esempio:\\
        \texttt{42["message", \{"sender":"Mario", "message":"Ciao a tutti!"\}]}
        \item Subito dopo, si vedranno pacchetti WebSocket simili in direzione opposta (Server -> Client) che rappresentano il "broadcasting" dell'evento \texttt{UploadChat} verso gli altri client connessi per aggiornare le loro interfacce.
    \end{itemize}
\end{enumerate}

\textbf{Conclusione:}\\
\textbf{La cattura dimostra visivamente il più grande vantaggio dei WebSocket:} al contrario del polling \texttt{HTTP}, non c'è bisogno di instaurare continuamente nuove connessioni o inviare pesanti Header \texttt{HTTP} per ogni singolo messaggio. Una volta effettuato l'handshake e l'Upgrade (101 Switching Protocols), il tunnel \texttt{TCP} resta permanentemente aperto e i messaggi viaggiano al suo interno come frame leggerissimi e a bassissima latenza.
\end{tcolorbox}

\end{document}
